\chapter{Mười Lát Cắt Học Sinh Và Màn Hình}
\markboth{Mười Lát Cắt Học Sinh Và Màn Hình}{Mười Lát Cắt Học Sinh Và Màn Hình}

\latcat{1. Vào Lớp Là Mở Máy}{Nhiều học sinh bước vào lớp mà tay đã tìm điện thoại trước khi tìm tập.\\
\textit{(Many students enter the classroom with their hand reaching for the phone before reaching for their notebook.)}\\
Phản xạ đó nói một điều rất buồn: học chưa bắt đầu mà sự phân tâm đã vào chỗ trước.\\
\textit{(That reflex says something sad: learning has not even begun, yet distraction has already taken its seat.)}}{A bad class often starts before the teacher says the first word.}

\latcat{2. Giờ Ra Chơi Không Có Tiếng Cười}{Ngày xưa ra chơi nghe tiếng người.\\
\textit{(Recess used to sound like human voices.)}\\
Bây giờ nhiều góc sân đầy học sinh mà im kiểu lạnh, vì ai cũng đang ở trong một cái hố nhỏ của màn hình.\\
\textit{(Now many corners of the yard are full of students yet carry a cold silence, because everyone is trapped in a small pit of screen-light.)}}{Crowds do not guarantee community.}

\latcat{3. Học Bài Mười Phút, Lướt Ba Mươi Phút}{Một số em rất thích cảm giác mình đã ngồi vào bàn học.\\
\textit{(Some students love the feeling that they have sat down to study.)}\\
Nhưng đời không tính bằng tư thế ngồi. Nó tính bằng khoảng tập trung thật, và điện thoại thường là kẻ cướp sạch khoảng đó.\\
\textit{(But life does not measure by sitting posture. It measures real concentration, and the phone is often the thief that steals it.)}}{Pretending to study is one of the easiest lies to tell yourself.}

\latcat{4. Toilet Cũng Không Tha}{Có em vào toilet không để đi vệ sinh cho xong mà để có cớ lướt thêm một chút.\\
\textit{(Some students go to the toilet not just to use it, but to gain an excuse for more scrolling.)}\\
Đó không còn là giải trí nữa. Đó là dấu hiệu một người không còn giữ nổi mình ở những chỗ nhỏ nhất.\\
\textit{(That is no longer simple entertainment. It is a sign of a person who cannot hold themselves together even in the smallest places.)}}{Lack of control in tiny spaces often exposes deeper weakness.}

\latcat{5. Đêm Mất Vì Ánh Sáng Xanh}{Nhiều học sinh than mệt, nhưng trước khi ngủ vẫn nằm dưới chăn thêm một giờ với màn hình.\\
\textit{(Many students complain of exhaustion, yet still spend another hour under the blanket with a screen before sleep.)}\\
Rồi sáng hôm sau lại gọi hậu quả đó là áp lực bí ẩn.\\
\textit{(The next morning they call the result a mysterious pressure.)}}{Not every suffering is mysterious. Some suffering is just repeated stupidity.}

\latcat{6. Vừa Đi Vừa Cắm Mặt}{Đi bộ ngoài đường mà mắt không nhìn đường là một kiểu coi thường thực tại.\\
\textit{(Walking on the street without watching the road is a way of disrespecting reality.)}\\
Con người đang ở giữa đời thật mà lại sống như thể mình chỉ tạm ghé qua cho xong.\\
\textit{(A person standing in real life is living as if they are only passing through it.)}}{A bent neck can slowly become a bent way of living.}

\latcat{7. Mọi Cảm Xúc Đều Phải Có Nhạc Nền}{Buồn cũng bật điện thoại.\\
\textit{(Sad, and the phone comes out.)}\\
Chán cũng bật điện thoại.\\
\textit{(Bored, and the phone comes out.)}\\
Rảnh cũng bật điện thoại.\\
\textit{(Free, and the phone comes out.)}\\
Lâu dần, học sinh không còn biết đi qua cảm xúc bằng nội lực của mình nữa.\\
\textit{(Over time, students stop knowing how to pass through emotion using their own inner strength.)}}{Constant soothing weakens emotional muscle.}

\latcat{8. Không Chịu Nổi Một Trang Sách}{Một video mười lăm giây xem được.\\
\textit{(A fifteen-second video feels easy.)}\\
Nhưng một trang sách dài một chút đã thấy ngứa tay, ngứa mắt, ngứa đầu.\\
\textit{(But one slightly longer page already makes the hands, eyes, and mind itch.)}}{The brain trained on fragments becomes allergic to depth.}

\latcat{9. Bị Chụp Đứt Mọi Kỷ Niệm}{Đi chơi với bạn mà đầu óc bận nghĩ góc chụp hơn là cuộc gặp là một kiểu mất mát rất hiện đại.\\
\textit{(Going out with friends while caring more about the shot than the meeting is a very modern kind of loss.)}\\
Kỷ niệm bị thay bằng tư liệu.\\
\textit{(Memory gets replaced by content material.)}}{Not everything worth keeping needs to become a post.}

\latcat{10. Sự Trì Hoãn Có Ống Kính}{Điện thoại cho sự trì hoãn một bộ mặt rất đẹp.\\
\textit{(The phone gives procrastination a very attractive face.)}\\
Người ta không còn thấy mình đang trốn việc. Người ta chỉ thấy mình đang xem một chút, nghỉ một chút, cập nhật một chút.\\
\textit{(People no longer see themselves escaping work. They only see themselves watching a little, resting a little, updating a little.)}}{The most dangerous delays are the ones that look harmless.}

\ngancach